\providecommand{\todo}[1]{\textbf{[TODO: #1]}}

บทนี้สรุปผลการดำเนินงานตามวัตถุประสงค์ พร้อมข้อจำกัดและข้อเสนอแนะสำหรับการพัฒนาต่อ

\section{สรุปผล}

\begin{enumerate}
    \item \textbf{วัตถุประสงค์ข้อ 1:} พัฒนาระบบถอดความเสียงบรรยายผลการตรวจชิ้นเนื้อที่ประมวลผลบนเครื่องส่วนบุคคลได้ในรูปแบบเว็บแอปพลิเคชันที่ให้บริการบนเครื่องผู้ใช้ ผ่านการทดสอบการทำงาน 9 จาก 9 โมดูล และตรวจว่าไฟล์ CSS, JS และฟอนต์ฝั่งหน้าเว็บเป็นไฟล์ในเครื่อง ผู้จัดทำรายงานว่าได้ทดสอบโดยปิดการเชื่อมต่อเครือข่ายทั้งหมดแล้วระบบยังทำงานได้ \todo{แนบหลักฐานการทดสอบตัดเครือข่าย (วันที่ วิธีที่ใช้ เช่น ถอดสาย LAN และปิด Wi-Fi และภาพหน้าจอหรือบันทึกผล)}
    \item \textbf{วัตถุประสงค์ข้อ 2:} ระบบที่ประกอบด้วย afftdn, INT8, ข้อความนำ CAP และ Beam Size 1 ลด WER เฉลี่ยจาก 36.27\% เป็น 29.05\% (ผลต่าง 7.21 จุด ช่วงความเชื่อมั่น 95\% 6.26--8.14) และจาก 7.25\% เป็น 2.78\% เมื่อปรับรูปแบบมิติ ในหมวดเสียงตู้ดูดควันลดจาก 44.88\% เป็น 31.94\% การทดลองแยกองค์ประกอบบน 5 กรณีไม่พบว่า afftdn เพียงอย่างเดียวลด WER จึงยังไม่มีหลักฐานว่าการกรองเสียงมีส่วนต่อการลดนี้
    \item \textbf{วัตถุประสงค์ข้อ 3:} ตัวสกัดแบบกำหนดแน่นอนกรอกฟิลด์และสร้างรายงาน PDF ได้ สัดส่วนเคสที่ฟิลด์วิกฤตอย่างน้อยหนึ่งฟิลด์ผิดลดจาก 27.5\% เป็น 19.6\% (ผลต่างแบบจับคู่ 7.9 จุด ช่วงความเชื่อมั่น 5.3--10.6) แต่ร้อยละ 10.0 ของเคสผิดจากตัวสกัดเองแม้ใช้ข้อความต้นฉบับ (หมวดที่ผู้พูดสลับลำดับ) Macro-F1 ของ 11 ฟิลด์เพิ่มเล็กน้อยจาก 96.26\% เป็น 96.97\% การเว้นว่างลดจาก 332 เหลือ 196 ช่อง แต่ค่าที่กรอกผิดเพิ่มจาก 146 เป็น 164 ช่อง ตัวปรับข้อความผ่านการทดสอบหน่วย 16 จาก 16 กรณี
    \item \textbf{วัตถุประสงค์ข้อ 4:} เวลาประมวลผลเฉลี่ยลดจาก 20.04 เป็น 11.50 วินาทีต่อกรณี (เร็วขึ้น 1.74 เท่า มัธยฐาน 16.30 เป็น 9.02 วินาที) การลด WER และเวลามีนัยสำคัญทางสถิติเมื่อจับคู่ระดับกรณี แต่เป็นการเปรียบเทียบทั้งระบบและกรณีทดสอบไม่เป็นอิสระต่อกันอย่างสมบูรณ์
\end{enumerate}

ผลทั้งหมดได้จากชุดข้อมูลเสียงสังเคราะห์ 1,000 กรณี ซึ่งซ้อนทับเสียงพัดลมตู้ดูดควันที่บันทึกจริง ยังไม่ผ่านการทดสอบทางคลินิกในห้องปฏิบัติการจริง และระบบเป็นเครื่องมือช่วยบันทึกที่ผู้ใช้ต้องตรวจทานก่อนยืนยันผล

\section{ข้อจำกัดที่สำคัญ}
\label{sec:limits}

\begin{itemize}
    \item \textbf{ข้อมูลทดสอบ:} เสียงสังเคราะห์ 2 เสียง สำเนียงอเมริกัน หมวดพูดเร็วได้จากการเร่งความเร็วเสียง ข้อความ เฉลย และข้อความนำมาจากแม่แบบใกล้เคียงกัน กรณีทดสอบจึงไม่เป็นอิสระต่อกัน
    \item \textbf{เฉลย:} สร้างจากสคริปต์ ไม่ได้ผ่านการตรวจรายกรณีโดยพยาธิแพทย์ ถูกแก้ 3 จุดภายหลังการประเมินรอบแรก และ 4 ฟิลด์ไม่มีค่าบวก จึงไม่ได้ทดสอบความสามารถในการตรวจจับ
    \item \textbf{การปรับตัวสกัดบนชุดทดสอบ:} ตัวสกัดถูกปรับหลังเห็นข้อผิดพลาดบนชุดเดียวกัน (Macro-F1 ก่อนปรับ 94.47\%) และเกณฑ์ WER B ถูกกำหนดหลังเห็นผลของเกณฑ์ A ตัวเลขการสกัดจึงเป็นขอบเขตบน
    \item \textbf{จุดอ่อนของตัวสกัด:} ขนาดสิ่งส่งตรวจและขนาดก้อนผิดทั้งหมวดที่ผู้พูดสลับลำดับ ก้อนลุกลามถูกสกัดผิดใน 12 จาก 100 กรณีของหมวด 6 และระยะห่างขอบตัดด้านลึกแย่กว่า baseline (ถอดคำว่า deep ผิดใน 41 กรณีที่ baseline ถูก)
    \item \textbf{ค่าที่กรอกผิดและธงเตือน:} ค่าที่กรอกผิดเพิ่มจาก 146 เป็น 164 ช่อง (ฟิลด์วิกฤตจาก 116 เป็น 137) และธงเตือนครอบคลุมเคสวิกฤตผิดร้อยละ 65.8 ผู้ใช้จึงต้องตรวจทานทุกฟิลด์วิกฤต
    \item \textbf{กฎของตัวปรับข้อความ:} กฎที่เปลี่ยนข้อมูลอัตโนมัติ (เช่น แทนค่าเดิมด้วยค่าหลังเมื่อพบ no ระหว่างค่าที่มีหน่วย และแทน tumor ด้วย mass) ผ่านชุดทดสอบ 16 กรณีที่เขียนหลังพบข้อบกพร่อง แต่ยังไม่ได้ยืนยันกับถ้อยคำจริงหรือโดยพยาธิแพทย์
    \item \textbf{การเปรียบเทียบและ Ablation:} baseline ต่างจาก PathoWhisper หลายองค์ประกอบพร้อมกัน Ablation ทำเพียง 5 กรณีและเปลี่ยนไลบรารีพร้อม INT8 การเปรียบเทียบกับ Vosk และ wav2vec 2.0 ทำบน 20 กรณีด้วยตัวแปลงที่ผู้จัดทำเขียนเอง และยังไม่ได้ทดสอบ Whisper รุ่นใหญ่ขึ้น
    \item \textbf{ระบบเว็บ:} ทดสอบเฉพาะเส้นทางหลัก ฐานข้อมูลสองที่ยังไม่มีกลไกกระทบยอด และใช้เซิร์ฟเวอร์สำหรับพัฒนากับใบรับรองที่ออกเอง
    \item \textbf{การตรวจสอบกับการใช้งานจริง:} ยังไม่ได้ทดสอบกับเสียงมนุษย์จริงหรือกับพยาธิแพทย์ในห้องปฏิบัติการ
\end{itemize}

\section{ข้อเสนอแนะ}

\subsection{ด้านการประเมินและตัวสกัด}

\begin{itemize}
    \item \textbf{ชุดทดสอบอิสระ:} ตรึงตัวสกัดและสร้างชุดทดสอบใหม่ (เช่น 100--200 กรณี) ด้วยแม่แบบและถ้อยคำต่างจากชุดปัจจุบัน แล้วรายงานผลชุดนั้นเป็นผลหลัก พร้อมวิเคราะห์ผลตามผู้พูดหรือแม่แบบเพื่อไม่ให้ประเมินความมั่นใจเกินจริง
    \item \textbf{แก้การสลับลำดับ:} ปรับตัวสกัดให้จับค่าตามบริบทของถ้อยคำ (เช่น คำว่า specimen, mass) แทนการพึ่งลำดับที่พูด หรือพิจารณาการสกัดเชิงความหมาย (Clinical NLP) ที่ตรวจสอบย้อนกลับได้ และตรวจผลรายกรณีของหมวด 2 และหมวด 6
    \item \textbf{ธงเตือนระดับฟิลด์:} เพิ่มธงเตือนที่ชี้ฟิลด์ที่ผิดเมื่อพบการแก้ไขกลางประโยค ขอบตัดที่ไม่มีคำว่า deep หรือผลถอดที่ขัดกับฟิลด์อื่น และวัด Flag Coverage ระดับฟิลด์
    \item \textbf{ขอบตัดด้านลึก:} สืบหาสาเหตุที่ Recall ลดลงเมื่อเทียบกับ baseline และขยายการสกัดไปยังขอบตัดด้านอื่น
    \item \textbf{ด้านของเต้านม:} ทดสอบอคติของข้อความนำด้วยข้อความนำที่เป็นกลางหรือสลับลำดับ right และ left และรายงานเมทริกซ์ความสับสนซ้ายขวา
    \item \textbf{Ablation ที่ทำซ้ำได้:} รันบนทั้งหมวดเสียงตู้ดูดควัน 100 กรณี ตั้ง temperature เป็น 0 และ seed คงที่ เก็บผลรายกรณี แยกขั้น silenceremove และใช้ faster-whisper INT8 ที่ประมวลผลเสียงเหมือนกันเป็นตัวควบคุม
    \item \textbf{กฎของตัวปรับข้อความ:} ให้พยาธิแพทย์ตรวจว่ากฎแก้ไขกลางประโยค (โดยเฉพาะ no ระหว่างค่าที่มีหน่วย) ไม่เปลี่ยนความหมายทางคลินิก และขยายชุดทดสอบด้วยกรณีที่เขียนโดยผู้ไม่ได้พัฒนากฎ
    \item \textbf{การตรวจสอบกับการใช้งานจริง:} ทดสอบกับเสียงมนุษย์จริง โดยเฉพาะผู้พูดชาวไทยที่พูดอังกฤษทางการแพทย์ ในห้องที่มีเสียงตู้ดูดควันจริง ภายใต้การรับรองด้านจริยธรรม และให้พยาธิแพทย์ตรวจถ้อยคำเฉลย
    \item \textbf{การเปรียบเทียบอย่างเป็นธรรม:} ขยายการเปรียบเทียบ Vosk และ wav2vec 2.0 พร้อมการแปลงคำอ่านตัวเลข จาก 20 กรณีเป็นชุดเต็มหรือชุดสุ่มที่ใหญ่ขึ้น ตรวจตัวแปลงโดยผู้ที่ไม่ได้เขียน และเพิ่มแบบจำลองที่แข็งกว่า เช่น Whisper รุ่นใหญ่ขึ้น พร้อมวัดเวลาบนเครื่องจริง
    \item \textbf{การ Fine-tune:} ทดลอง Fine-tune หรือ LoRA กับเสียงจริงที่ได้รับอนุญาต แล้วเปรียบเทียบกับแนวทาง Prompt
\end{itemize}

\subsection{ข้อกำหนดเชิงระบบก่อนนำไปใช้ในโรงพยาบาล}

ข้อเสนอต่อไปนี้เป็นข้อกำหนดที่ควรมีก่อนใช้งานจริง ระบบต้นแบบปัจจุบันยังไม่ได้ทำครบทุกข้อ

\begin{itemize}
    \item \textbf{การลงนามโดยพยาธิแพทย์:} ผลที่ระบบกรอกเป็นเพียงฉบับร่าง ต้องบังคับให้พยาธิแพทย์ตรวจทานและยืนยันก่อนลงนาม และรายงาน PDF และ DOCX ก่อนยืนยันควรมีลายน้ำระบุว่าเป็นฉบับร่างที่รอการตรวจยืนยัน
    \item \textbf{การควบคุมสิทธิ์ตามบทบาท (RBAC):} แยกสิทธิ์ของผู้ร่างบันทึก (ผู้ช่วยหรือแพทย์ประจำบ้าน) ออกจากพยาธิแพทย์ผู้รับรองผลสุดท้าย ให้ผู้ดูแลระบบเป็นผู้สร้างบัญชีแทนการเปิดสมัครเอง และกำหนดอายุของเซสชัน
    \item \textbf{ร่องรอยการตรวจสอบ (Audit Trail):} เก็บเสียงต้นฉบับ ข้อความถอดความดิบ ค่าที่ระบบกรอก และประวัติการแก้ไขทุกครั้ง (ผู้แก้ เวลา ค่าเดิม ค่าใหม่) ในรูปแบบที่แก้ไขย้อนหลังไม่ได้ และกำหนดระยะเวลาเก็บและการลบเสียงให้สอดคล้องกับ พ.ร.บ. คุ้มครองข้อมูลส่วนบุคคล พ.ศ. 2562 โดยเข้ารหัสข้อมูลที่จัดเก็บ
    \item \textbf{การช่วยตรวจทาน:} ไฮไลต์ฟิลด์ที่ถูกแก้ไขกลางประโยคหรือมีความเชื่อมั่นต่ำ แจ้งเตือนเมื่อระบบลบหรือแทนค่าเดิม แสดงข้อความถอดความดิบเทียบกับค่าที่สกัด และเปิดฟังเสียงเทียบได้ เพื่อลดความเสี่ยงที่ผู้ใช้เชื่อผลอัตโนมัติเกินควร
    \item \textbf{ฐานข้อมูล:} กำหนดฐานข้อมูลหลักเพียงแห่งเดียวพร้อมกลไกกระทบยอดกับฐานสำรอง เพื่อไม่ให้ข้อมูลซ้ำหรือหาย
    \item \textbf{การใช้งานจริง:} เปลี่ยนเซิร์ฟเวอร์สำหรับพัฒนาและใบรับรองที่ออกเองเป็นระดับใช้งานจริง และทดสอบความปลอดภัย การใช้งานพร้อมกัน และการแยกสิทธิ์ ระบบนี้ยังไม่ใช่เครื่องมือแพทย์ที่ได้รับการรับรอง
    \item \textbf{การขยายและเชื่อมต่อ:} ขยายไปมะเร็งลำไส้ใหญ่และต่อมไทรอยด์ตามโปรโตคอล CAP ใช้แป้นเหยียบ USB สั่งบันทึกเสียง และเชื่อมต่อ LIS ผ่าน HL7 หรือ FHIR
\end{itemize}

\section{บทส่งท้าย}

โครงงานนี้แสดงให้เห็นว่าระบบแปลงเสียงเป็นข้อความแบบออฟไลน์ที่ปรับให้เหมาะกับเสียงรบกวนและศัพท์เฉพาะของห้องตรวจชิ้นเนื้อ สามารถลดความผิดพลาดของการถอดความและเวลารอเมื่อเทียบกับ Whisper Small ได้ในชุดทดสอบสังเคราะห์ แต่การสกัดข้อมูลลงฟิลด์ยังมีจุดอ่อนที่เกี่ยวข้องกับข้อมูลสำคัญทางคลินิก และตัวเลขที่ได้ยังไม่ผ่านการยืนยันด้วยข้อมูลอิสระ หากแก้ข้อบกพร่องที่พบและผ่านการทดสอบกับเสียงจริง ระบบอาจช่วยลดการจดซ้ำและการสัมผัสแป้นพิมพ์ขณะปฏิบัติงาน โดยยังคงให้ผู้เชี่ยวชาญเป็นผู้ตัดสินใจขั้นสุดท้าย
